\section[Conclusão]{Conclusão}

Participar do projeto Estudify como \ac{ages} IV foi, para mim, o ponto alto da minha trajetória dentro da \ac{ages}. Depois de passar pelos papéis de desenvolvedor, analista de dados e arquiteto ao longo dos semestres anteriores, encerrar esse ciclo na posição de maior senioridade, responsável pela gestão de uma equipe inteira, deu um sentido de conclusão à minha caminhada na Agência.

Sempre tive um perfil mais técnico, e assumir a gestão me tirou da zona de conforto. Precisei deslocar minha atenção do código para as pessoas: planejar as \textit{sprints}, escrever e revisar as \ac{us}, organizar o time em \textit{squads} e, acima de tudo, manter cada integrante avançando. Foi justamente nessa frente humana que encontrei meus maiores desafios. Reconheço que, em alguns momentos, demorei a agir quando uma tarefa começava a atrasar e, sem o acompanhamento devido, esse pequeno atraso acabava se transformando em um risco concreto para a entrega da \textit{sprint}. Aprendi, na prática, que liderar exige uma postura ativa: não basta perguntar se está tudo bem, é preciso acompanhar de perto e se colocar na posição de cada um, considerando seu nível de conhecimento e suas dificuldades.

Ainda assim, meu perfil técnico continuou sendo útil ao longo do projeto. Em momentos que foi necessário, voltei ao código para auxiliar o time. Equilibrar essa atuação técnica com as responsabilidades de gestão foi um exercício constante, mas me mostrou que as duas dimensões podem se complementar.

Quanto à relação entre as disciplinas cursadas e a \ac{ages}, a que mais me preparou para esse papel foi Gerenciamento de Projetos de Software, cuja dependência fica evidente na prática: planejamento, definição de escopo, gerenciamento de riscos e acompanhamento de cronograma foram conceitos que apliquei diretamente ao estruturar o escopo em épicos, ao detalhar as histórias com critérios de aceite claros e ao organizar o time em \textit{squads}, com um \ac{ages} IV responsável por cada uma.

Sem dúvida, o Estudify foi o projeto mais marcante e completo de que participei. Pude contar com uma equipe genuinamente engajada. Foi também a entrega mais robusta da minha trajetória: além de cumprir todo o escopo previsto, entregamos funcionalidades adicionais como \textit{overdelivery}, disponibilizamos o aplicativo às \textit{stakeholders} por meio do TestFlight e construímos uma base real de questões do ENEM. Ver as clientes recebendo muito mais do que haviam pedido, me deu a certeza de que construímos algo com potencial real de uso.

Levo dessa experiência aprendizados que vão muito além do técnico. Como pontos positivos, destaco a importância de escrever tarefas bem detalhadas, prática que reduziu sensivelmente o retrabalho dos \ac{ages} menores, e a eficácia do modelo de \textit{squads} acompanhadas individualmente. Como pontos a melhorar, reconheço a necessidade de alinhar mais cedo a escrita das histórias com a modelagem do banco e de acompanhar o time de forma mais próxima e proativa nos momentos críticos. Mais do que qualquer entrega, foi nesta passagem que percebi o quanto a gestão de pessoas é desafiadora. Encerro o ciclo da \ac{ages} com a sensação de dever cumprido, tendo evoluído de um perfil puramente técnico para alguém capaz de liderar, comunicar e extrair o melhor de cada integrante de uma equipe, competências que pretendo levar comigo para o início da minha carreira como engenheiro de software.
